\scanpage{77}
\unsure{Pour|P.} $i$ grand $G,G'$ proviennent de $G_i,G'_i$ \inserted{sch. en \unsure{grpes|gr.} comm de t.f.} plats sur $X_i$
\deleted{et $\xi$ d'une corr.}
$G_i|U_i=A_i$ et $G'_i|\unsure{U_i|U_i'|U_i^*}=A'_i$ des sch. abéliens,
$\xi$ provient d'un élém de corr. div. $\xi'$ sur $A_i\times A'_i$,
$T$ provient d'un tore $T_i\subset G'_{iY_i}$, $S$ provient d'un sous-schéma plat $S_i\subset G_{iY_i}$.
De plus, \deleted{\unsure{dans le cas}} \deleted{\unsure{Encore}} on peut \deleted{\unsure{étant|sans|était} \unsure{conditions|restriction}} supposer $S$ connexe \inserted{\unsure{\uline{lisse}|$G_{\underline m}$}} (sinon le remplacer par $S^0$) et
\deleted{\unsure{de la partie unipotente} \unsure{Comme} si $p=\operatorname{car}k>0$ Comme le \unsure{remplacer|quotient|montre} par $p^nS$}
extension \deleted{d'une} successive d'une V.A, d'un tore \deleted{\unsure{pour} $S$ \unsure{propre}} \unsure{et|à} \unsure{de|si} \deleted{\unsure{un grpe}} grpe $G_a$.
On peut donc supposer le $=$ dévissage valable pour $S_i$.
\deleted{\unsure{Cela implique}}
On peut supposer enfin (remplaçant au besoin les $A_i$ par \deleted{$\unsure{A_i|A}$} $A_i[1/\ell]$) que $\ell$ inv. sur les $X_i$ \unsure{donc|sur} les $Y_i$.
Mais les $ {}_{\ell^{\unsure{\nu|n|d}}}S_i$ sont encore fini sur $Y_i$.
Bref, ttes les données proviennent de \deleted{\unsure{un}} données analogues, avec un $\unsure{\deleted{X_0}|\deleted{X}S_i|\widetilde X_i}$
$X_0$ \inserted{$=\operatorname{Spec}A_0$} \uline{de type fini} sur $\operatorname{Spec}\mathbf Z$.
De plus, si $\mathfrak p=\mathfrak m\cap A_0$, alors $A_0\mathbin{-}V$ induit $\widehat A_0\longrightarrow V$, \deleted{\unsure{comme}}
Donc, on peut remplacer les mots ``$A_0$ de type fini sur $\operatorname{Spec}\mathbf Z$'' par ``$A_0$ $\mathfrak p$-adique séparé et \unsure{complet|compl.}, avec $A_0/\mathfrak p$ de t.f. sur $\unsure{\operatorname{Spec}\mathbf Z|\deleted{\mathbf F_p}\mathbf Z}$, et $A_0\mathbin{-}V$ continu pour les topologi\unsure{es} adiques''.
On est ramené alors au \unsure{th.|thm.} pour cet $X_0,Y_0,G_0,G'_0,\xi_0,S_0,T_0$, et il relève du cas 1.
\hfill CQFD.
\margin{76}
